\chapter{多维随机变量及其分布}\label{ch:03}
\chaptermap{联合分布\\支撑与归一}{边缘与条件\\独立性判定}{二维变换\\和、比、雅可比}{次序与正态\\结构化计算}
\begin{examtip}
\textbf{考研核心：}联合密度、边缘密度、条件密度、独立性及和的分布。
\textbf{竞赛视角：}顺序区域、和比变换、正态旋转。
\textbf{测度关联：}积分上下限由支撑的截面决定，分布变换是概率质量的搬运。
\end{examtip}
\section{联合分布、边缘分布与支撑}\label{sec:31}
\src{茆诗松, §3.1—3.2；李东风, §3.1—3.3}
\begin{definition}{联合分布}
$F_{X,Y}(x,y)=\Pp(X\le x,Y\le y)$。
联合密度$f$满足$f\ge0$且$\iint f=1$，事件概率为事件域上的积分。
$f_X(x)=\int f(x,y)\dd y$，$f_Y(y)=\int f(x,y)\dd x$；离散情形沿行列求和。
联合CDF的矩形增量必须非负，仅逐坐标非减不足以成为联合CDF。
\end{definition}
\begin{method}{画支撑、取截面、再积分}
适用：三角形、曲边区域上的密度。先写支撑不等式，固定$x$求$y$区间，固定$y$求$x$区间。
归一化常数、边缘密度和事件概率共用这张图；不同积分次序的上下限通常不同。
\end{method}
\worked{三角形上的均匀分布}{基础例题}{\star\star}{3-1}
{$f(x,y)=c$在$0<x<y<1$上，求$c$、边缘密度及$\Pp(X+Y<1)$。}
{先把概率看成面积比。}
{归一化给$c(\frac{1}{2})=1$，所以$c=2$。边缘为$f_X(x)=2(1-x)$、$f_Y(y)=2y$（$0<x,y<1$）。
方法一：$\Pp(X+Y<1)=\int_0^{\frac{1}{2}}\int_x^{1-x}2\dd y\dd x=\frac{1}{2}$。
方法二：映射$(x,y)\mapsto(1-y,1-x)$保持支撑和面积，并交换$X+Y<1$与$X+Y>1$；边界面积零，故各为$\frac{1}{2}$。}
{积分可计算一般区域；对称法在本题直接平分质量。}
{支撑非矩形，不能把$X,Y$当独立均匀。}
\worked{非均匀三角密度}{考研风格}{\star\star}{3-2}
{$f(x,y)=kx$在$0<y<x<1$上，求$k$及$f_Y$。}
{沿竖截面求常数，沿横截面求边缘。}
{$1=\int_0^1\int_0^xkx\dd y\dd x=\frac{k}{3}$，故$k=3$。
固定$y$时$y<x<1$，故$f_Y(y)=\int_y^13x\dd x=\frac32(1-y^2)$，$0<y<1$，其他为0。
检查$\int_0^1\frac32(1-y^2)\dd y=1$。}
{画区域能防止把两个边缘使用同一上下限。}
{密度中有$x$不代表$Y$均匀，积分区间仍依赖$y$。}
\worked{均匀点的距离带}{竞赛风格}{\star\star}{3-3}
{独立$X,Y\sim U(0,1)$，求$D=|X-Y|$密度。}
{单位正方形中带状区域的补集是两三角形。}
{$0\le d\le1$时$F_D(d)=1-(1-d)^2$，故$f_D(d)=2(1-d)$。
另一种看法是$X-Y$密度为$1-|t|$（$|t|<1$），绝对值合并$t=\pm d$两支，得到同式。
归一化检查$\int_0^12(1-d)\dd d=1$。}
{几何补域与卷积互相验证。}
{绝对值变换合并两支，但$d=0$没有原子。}
\begin{figure}[htbp]\centering
\begin{tikzpicture}[scale=4]
\fill[Teal!12](0,0)--(0,1)--(1,1)--cycle;
\draw[->,Ink](0,0)--(1.1,0)node[right]{$x$};\draw[->,Ink](0,0)--(0,1.1)node[above]{$y$};
\draw[Teal,thick](0,0)--(1,1)--(0,1);\draw[Gold,thick](0,1)--(1,0);
\draw[Rust,dashed](.3,.3)--(.3,1);\node[right]at(.3,.85){\small 竖截面};
\node at(.25,.55){$x<y$};\node[below]at(1,0){1};\node[left]at(0,1){1};
\end{tikzpicture}\caption{先看支撑三角形，再用$x+y=1$切分概率}\label{fig:triangle}\end{figure}

\section{条件分布与独立性}\label{sec:32}
\src{茆诗松, §3.2、§3.5；李东风, §3.6}
\begin{definition}{条件密度与独立}
离散情形记$p_{X,Y}(x,y)=\Pp(X=x,Y=y)$，$p_X(x)=\sum_y p_{X,Y}(x,y)$。当$p_X(x)>0$时，$\Pp(Y=y\mid X=x)=\frac{p_{X,Y}(x,y)}{p_X(x)}$。
有联合密度时，对$0<f_X(x)<\infty$可取$f_{Y\mid X}(y\mid x)=\frac{f_{X,Y}(x,y)}{f_X(x)}$。
$X,Y$独立当且仅当每个$x,y$都有$F_{X,Y}(x,y)=F_X(x)F_Y(y)$；离散时等价于所有取值的$p_{X,Y}=p_Xp_Y$，有密度时等价于$f_{X,Y}=f_Xf_Y$几乎处处。
条件密度只在边缘分布几乎处处的$x$上确定。
\end{definition}
\begin{method}{先核对支撑是否能分离}
适用：独立性判断、条件密度。若边缘密度同时为正的矩形内存在联合密度为零的正面积区域，则不独立。
公式的代数可分离还不够，支撑也需以几乎处处意义相容。
\end{method}
\worked{给定上边界后的均匀性}{考研风格}{\star\star}{3-4}
{对例\ref{ex:3-1}，求$X\given Y=y$及$\Pp(X<\frac{Y}{2})$。}
{固定$y$后，横截面是$(0,y)$。}
{方法一：$f_{X\mid Y}(x\mid y)=\frac{2}{2y}=\frac{1}{y}$，所以条件概率$\Pp(X<\frac{y}{2}\given Y=y)=\frac{1}{2}$，再全概率积分得$\frac{1}{2}$。
方法二：直接$\int_0^1\int_0^{\frac{y}{2}}2\dd x\dd y=\frac{1}{2}$。
条件密度同时说明$X\given Y=y\sim U(0,y)$，其范围随$y$变，故不独立。}
{先条件化可把二维区域变成长度比例。}
{条件密度的分母是固定条件变量的边缘密度。}
\worked{可分离密度}{基础例题}{\star}{3-5}
{$f(x,y)=4xy$在$0<x,y<1$上，判断独立并求$\Pp(X<\frac{1}{2},Y<\frac{1}{2})$。}
{边缘密度与支撑均可分离。}
{$f_X(x)=2x,f_Y(y)=2y$，乘积等于联合密度，所以独立。
每个边缘的左半区概率为$\int_0^{\frac{1}{2}}2x\dd x=\frac{1}{4}$，联合概率为$\frac{1}{16}$。
直接二重积分亦为$(\int_0^{\frac{1}{2}}2x\dd x)^2$。}
{乘积密度加矩形支撑提供直接判据。}
{若同一公式只定义在三角区域，结论可能改变。}
\worked{零概率条件的正确语言}{测度论选学}{\star\star}{3-6}
{独立连续$X,Y\sim U(0,1)$，如何解释“给定$X=x$时$Y$仍均匀”？}
{不把零概率事件当成普通条件概率分母。}
{联合密度为1，边缘$f_X(x)=1$，所以可选条件密度$f_{Y\mid X}(y\mid x)=\ind_{(0,1)}(y)$（$0<x<1$）。
在$x\notin(0,1)$处也可将条件密度取为同一个均匀密度。它满足对任意Borel集$A,B$，
\[\Pp(X\in A,Y\in B)=\int_A f_X(x)\left[\int_B f_{Y\mid X}(y\mid x)\dd y\right]\dd x,\]
其中$f_X(x)=\ind_{(0,1)}(x)$，因此外层实际只在$A\cap(0,1)$上积分。
这条积分恒等式才是条件分布的含义；单个$x$处改变版本不影响恒等式。}
{密度比是条件分布的一种版本，不是$\frac{0}{0}$。}
{条件分布在零边缘质量点上可以任意指定。}

\section{二维变换与雅可比}\label{sec:33}
\src{茆诗松, §3.3；李东风, §3.4}
\begin{theorem}{二维密度变换}
设$(X,Y)$有联合密度。除原分布概率为0的集合外，定义域可分成有限或可数个不交开分支，每支上的$T$为$C^1$一一映射且雅可比行列式非零，因而有$C^1$逆变换。则$(U,V)=T(X,Y)$的密度几乎处处为
\[f_{U,V}(u,v)=\sum_{(x,y)\in T^{-1}(u,v)}f_{X,Y}(x,y)\left|\det\dfrac{\partial(x,y)}{\partial(u,v)}\right|.\]
求和只对有效支撑内的逆像进行。公式由二维积分换元得到，绝对雅可比表示面积伸缩率。
\end{theorem}
\begin{method}{和与比例配成一对}
适用：正变量的和、比值、分配比例。设$S=X+Y$、$R=\frac{X}{X+Y}$，先反解$X=SR,Y=S(1-R)$，雅可比绝对值为$S$。
若只求一个变量，也可先条件化或做事件域积分验证。
\end{method}
\worked{指数变量的和与比例}{考研风格}{\star\star\star}{3-7}
{独立$X,Y\sim\Exp(\lambda)$，证明$R=\frac{X}{X+Y}$均匀且与$S=X+Y$独立。}
{变换把正象限映为$s>0,0<r<1$。}
{方法一：逆变换的导数矩阵为$\begin{pmatrix}r&s\\1-r&-s\end{pmatrix}$，行列式$-s$。
因此$f_{S,R}(s,r)=\lambda^2se^{-\lambda s}$，分解为Gamma密度与$U(0,1)$密度。
方法二：给定$S=s$时，$f_{X\mid S}(x\mid s)=\frac{1}{s}$（$0<x<s$）；故$\frac{X}{s}$在$(0,1)$均匀且条件分布不随$s$变化，得到独立。}
{变换法一次给出联合结构，条件法解释为何比例不受总量影响。}
{率不同时一般不独立，不能只看两个变量都叫指数分布。}
\worked{两个均匀变量的乘积}{考研风格}{\star\star}{3-8}
{独立$X,Y\sim U(0,1)$，求$Z=XY$密度。}
{固定一个变量，积分条件概率。}
{$0<z<1$时
\[F_Z(z)=\int_0^z1\dd x+\int_z^1\dfrac z x\dd x=z-z\log z,\qquad f_Z(z)=-\log z.\]
也可令$(Z,V)=(XY,X)$，逆变换$Y=\frac{Z}{V}$，雅可比$\frac{1}{v}$，范围$0<z<v<1$，边缘$\int_z^1v^{-1}\dd v=-\log z$。}
{条件法省去写完整二维变换；雅可比法更系统。}
{积分下限是$z$，因为$Y<1$要求$\frac{z}{v}<1$。}
\worked{极坐标与Rayleigh分布}{竞赛风格}{\star\star}{3-9}
{独立标准正态$X,Y$，求$R=\sqrt{X^2+Y^2}$和角度$\Theta$的联合密度。}
{圆对称密度适合极坐标。}
{取$x=r\cos\theta,y=r\sin\theta$，$r>0,0\le\theta<2\pi$，绝对雅可比$r$，故
$f_{R,\Theta}(r,\theta)=\frac{r}{2\pi}e^{-\frac{r^2}{2}}$。
分离后$f_R(r)=re^{-\frac{r^2}{2}}$、$\Theta\sim U(0,2\pi)$，且两者独立。
积分尾概率为$\Pp(R>r)=e^{-\frac{r^2}{2}}$。}
{旋转对称性先预测角度均匀，雅可比完成严格计算。}
{漏掉因子$r$会使半径在原点附近质量错误。}

\section{最大值、最小值与次序统计}\label{sec:34}
\src{李东风, §3.5；一般次序密度}
\begin{theorem}{iid次序统计量}
若$X_1,\ldots,X_n$独立同CDF为$F$，则$\Pp(X_{(n)}\le x)=F(x)^n$，$\Pp(X_{(1)}>x)=(1-F(x))^n$。
连续密度存在时
\[f_{X_{(k)}}(x)=\dfrac{n!}{(k-1)!(n-k)!}F(x)^{k-1}(1-F(x))^{n-k}f(x).\]
\end{theorem}
\begin{proof}
最大最小式由全部样本的交事件与独立性。
第$k$个值落在$(x,x+\dd x)$时，选择1个样本落在窄条，$k-1$个在左、$n-k$个在右；主阶系数为$n\binom{n-1}{k-1}$，两个以上落窄条只贡献高阶项，得到密度。
\end{proof}
\begin{method}{极值先用CDF，不先找联合密度}
适用：所有不超过或所有超过某阈值。一般次序量把“第$k$个不超过$x$”转成至少$k$个样本不超过$x$，即二项尾概率。
求极差时才通常需要最小值和最大值的联合密度。
\end{method}
\worked{两个均匀变量的最大值}{基础例题}{\star}{3-10}
{独立$X,Y\sim U(0,1)$，求$M=\max(X,Y)$密度。}
{同一事件可用乘法或面积。}
{方法一：$F_M(m)=\Pp(X\le m)\Pp(Y\le m)=m^2$（$0<m<1$），故$f_M(m)=2m$。
方法二：单位正方形内两个坐标均不超过$m$的区域是边长$m$的正方形，面积$m^2$。
若改最小值，CDF为$1-(1-m)^2$，密度$2(1-m)$。}
{图形与独立乘积表达的是同一矩形事件。}
{最大值相对单个均匀变量更集中于较大值，其CDF满足$m^2<m$（$0<m<1$）。这不等于统计学中的正偏态；本例密度$2m$的长尾朝左。}
\worked{均匀样本的极差}{考研风格}{\star\star\star}{3-11}
{$n\ge2$个iid $U(0,1)$样本，求极差$R=X_{(n)}-X_{(1)}$密度。}
{先求最小最大联合，再用差变换。}
{设最小$u$、最大$v$，选择两个端点样本有$n(n-1)$种，中间$n-2$个落在$(u,v)$，故$f_{u,v}=n(n-1)(v-u)^{n-2}$，$0<u<v<1$。
令$r=v-u$，雅可比为1，$0<u<1-r$，所以$f_R(r)=n(n-1)r^{n-2}(1-r)$（$0<r<1$）。}
{端点间的长度幂反映中间样本的自由放置。}
{不能把最大值与最小值当独立后相减卷积。}
\worked{均匀次序量的Beta结构}{考研风格}{\star\star}{3-12}
{5个iid $U(0,1)$样本，求中位数$X_{(3)}$密度及均值。}
{代入一般次序密度，再用对称性求均值。}
{$f(x)=\frac{5!}{2!2!}x^2(1-x)^2=30x^2(1-x)^2$，$0<x<1$。
密度关于$\frac{1}{2}$对称，故$\E X_{(3)}=\frac{1}{2}$。直接积分也可得$30(\frac{1}{4}-\frac{2}{5}+\frac{1}{6})=\frac{1}{2}$。
该分布为$\Beta(3,3)$。}
{Beta是均匀次序统计量的自然分布。}
{样本中位数不等于每个样本的平均值，虽在此有相同均值。}

\section{二维正态与线性变换}\label{sec:35}
\src{李东风, §5.1；统计笔记, §3}
\begin{definition}{二维正态}
均值$\mu_X,\mu_Y$、标准差$\sigma_X,\sigma_Y>0$、$|\rho|<1$的二维正态密度为
\[f(x,y)=\dfrac{\exp\{-\frac{z_x^2-2\rho z_xz_y+z_y^2}{2(1-\rho^2)}\}}{2\pi\sigma_X\sigma_Y\sqrt{1-\rho^2}},
\quad z_x=\dfrac{x-\mu_X}{\sigma_X},\ z_y=\dfrac{y-\mu_Y}{\sigma_Y}.\]
等价表示为$X=\mu_X+\sigma_XZ_1$、$Y=\mu_Y+\sigma_Y(\rho Z_1+\sqrt{1-\rho^2}Z_2)$，其中$Z_1,Z_2$独立标准正态。
\end{definition}
\begin{theorem}{正态条件分布与零相关}
二维正态的$Y\given X=x$为正态，均值$\mu_Y+\rho\frac{\sigma_Y}{\sigma_X}(x-\mu_X)$、方差$\sigma_Y^2(1-\rho^2)$。
联合正态中零协方差等价于独立。
\end{theorem}
\begin{proof}
由独立标准正态的线性表示，固定$X=x$便固定$Z_1$，剩下的$Z_2$仍独立标准正态，直接得到条件分布。
表示的雅可比绝对值为$\sigma_X\sigma_Y\sqrt{1-\rho^2}$，代回两个标准正态密度得到上式。
$\rho=0$时密度乘积分解；反向独立必有协方差零。
\end{proof}
\begin{method}{先构造独立正态方向}
适用：正态和差、条件分布。把相关变量写成独立标准正态的线性组合；计算均值协方差后判定。
“各自正态”不足以代替“联合正态”，零相关推出独立依赖后者。
\end{method}
\worked{正态和差的独立性}{考研风格}{\star\star}{3-13}
{独立$X,Y\sim\Normal(0,1)$，令$U=X+Y,V=X-Y$，求联合分布。}
{正交方向可用协方差或密度变换。}
{方法一：$(U,V)$联合正态，均值0，方差均为2，协方差$\Var X-\Var Y=0$，所以独立且各为$\Normal(0,2)$。
方法二：$x=\frac{u+v}{2},y=\frac{u-v}{2}$，绝对雅可比$\frac{1}{2}$，得到$f_{U,V}=\frac1{4\pi}e^{-\frac{u^2+v^2}{4}}$，恰好分解成两个$\Normal(0,2)$密度。}
{协方差法快，密度法展示独立的实际乘积结构。}
{若两变量方差不同，和与差一般相关。}
\worked{条件正态概率}{考研风格}{\star\star}{3-14}
{二维正态均值均0，方差均1，相关系数$\frac{1}{2}$，求$\Pp(Y>0\given X=2)$。}
{先确定条件均值和方差。}
{$Y\given X=2\sim\Normal(1,\frac{3}{4})$，所以
$\Pp(Y>0\given X=2)=1-\Phi(-\frac{2}{\sqrt3})=\Phi(\frac{2}{\sqrt3})$。
条件均值上移体现正相关；条件标准差为$\frac{\sqrt3}{2}$，小于无条件标准差1。}
{先写完整条件分布，再标准化不易漏平方根。}
{条件方差与$x$无关是联合正态的特殊结构。}
\worked{边缘正态却不联合正态}{概念辨析}{\star\star}{3-15}
{$Z\sim\Normal(0,1)$，独立符号$S$等概率取$\pm1$。令$X=Z,Y=SZ$，证明各自正态、不相关但不独立。}
{对称性决定边缘，绝对值关系揭示依赖。}
{$SZ$与$Z$同分布，所以两个边缘均标准正态。
$\E(XY)=\E(SZ^2)=\E S\E Z^2=0$，因而协方差零。
但$|X|=|Y|$几乎必然；若独立，连续变量$|X|,|Y|$相等的概率应为0，与此矛盾。
联合分布集中在两条对角线上，没有二维Lebesgue密度，但其协方差矩阵为 $I_2$，所以不能是联合正态。
还可直接检验线性组合：$X+Y$ 在 $S=-1$ 时等于零，在 $S=1$ 时等于 $2Z$，其分布是 $\tfrac12\delta_0+\tfrac12N(0,4)$。它在零处有概率质量 $1/2$，不可能是方差为二的正态分布；这违反联合正态的线性组合性质。}
{构造共同幅度与随机符号能保留边缘而改变依赖。}
{判定独立时，边缘形状与协方差都不足以替代联合信息。}
\begin{warning}
联合密度积分先看支撑。条件密度须对被条件变量取边缘。
正态和差独立要核对方差；次序统计量之间通常相关。雅可比应使用逆变换行列式的绝对值。
\end{warning}

\section*{图解、理解与巧解复盘}
\begin{insight}{边缘是截面质量，条件是截面内重新归一}
联合密度在整张区域上分配概率。沿$x$方向积掉得到$Y$边缘，固定$Y=y$后除以$f_Y(y)$把该截面的总质量缩放为1。
例\ref{ex:3-4}中截面长度为$y$，所以条件密度为$\frac{1}{y}$；这既解释条件均匀，又显示条件范围随$y$变化所产生的依赖。
\end{insight}
\begin{clevertrick}{正象限变成总量与比例矩形}
和比变换$X=SR,Y=S(1-R)$将$X,Y>0$变成$S>0,0<R<1$。
若两变量是独立同率Gamma$(a,\lambda)$与Gamma$(b,\lambda)$，代入并乘雅可比$S$后，联合核为
$s^{a+b-1}e^{-\lambda s}r^{a-1}(1-r)^{b-1}$。
利用$B(a,b)=\frac{\Gamma(a)\Gamma(b)}{\Gamma(a+b)}$归一化，得到$S\sim\Ga(a+b,\lambda)$、$R\sim\Beta(a,b)$且独立。
\end{clevertrick}
\begin{figure}[htbp]\centering
\begin{tikzpicture}[>=Stealth,scale=.85]
\draw[->,Ink](0,0)--(4,0)node[right]{$x$};\draw[->,Ink](0,0)--(0,3.7)node[above]{$y$};
\foreach \s in {1,2,3}{\draw[Teal](0,\s)--(\s,0);}
\draw[Gold,thick](0,0)--(3,3);\node at(2.8,2.6){\small 固定比例};\node at(1,-.7){固定总量：斜截线};
\draw[->,Ink,thick](4.5,1.5)--(5.5,1.5);
\draw[->,Ink](6,0)--(10,0)node[right]{$s$};\draw[->,Ink](6,0)--(6,3.7)node[above]{$r$};
\fill[Teal!9](6,0)rectangle(9.6,3);\draw[Teal](6,3)--(9.6,3);\node[left]at(6,3){1};
\draw[Gold,thick](6,1.5)--(9.6,1.5);\node at(8,2.3){$s>0,\ 0<r<1$};
\end{tikzpicture}\caption{和比坐标把斜截线与射线变为坐标线，简化支撑}\label{fig:sum-ratio}\end{figure}

\section*{章末分层习题}
\exercise{3-a}{基础}{$f(x,y)=cxy$在单位正方形上，求$c$及边缘密度。}
\exercise{3-b}{基础}{独立$U(0,1)$变量$X,Y$，求$\Pp(X+Y<\frac{1}{2})$。}
\exercise{3-c}{考研风格}{$f(x,y)=2$在$0<x<y<1$上，求$\E(X\given Y=y)$及$\Pp(Y>\frac{3}{4})$。}
\exercise{3-d}{考研风格}{独立$X\sim\Exp(2),Y\sim\Exp(3)$，求$\min(X,Y)$的分布。}
\exercise{3-e}{竞赛拓展}{独立$U(0,1)$变量$X,Y$，求$\frac{X}{Y}$的CDF。}
\exercise{3-f}{竞赛拓展}{独立标准正态$X,Y$，求$\Pp(X^2+Y^2>4)$。}
\exercise{3-g}{测度论选学}{令$Y=X$，$X\sim U(0,1)$，证明联合分布没有平面Lebesgue密度。}
\exercise{3-h}{测度论选学}{解释为什么修改条件密度在某个$x_0$处的值不改变联合概率，假设$X$绝对连续。}
\section*{章末方法速查}
\begin{quicktable}
截面积分&支撑已知&三角、曲边区域&上下限未分段\\
条件化&边缘密度为正&固定一维后简单&把零概率条件作分式\\
雅可比&可微逆变换&和比、极坐标&漏分支与支撑\\
正态线性结构&确为联合正态&和差、条件预测&只凭边缘正态推独立\\
\end{quicktable}
